\section{模型假设}

\begin{enumerate}[label=假设 \arabic*：]
  \item 巷道被视为由端点连接的三维线段。解释：附件给出了端点坐标和巷道端点编号，代码据此构建无向图 $G=(V,E)$。
  \item 突水点、矿工和出口均投影到最近巷道。解释：题面坐标可能不完全等于原始端点坐标，投影后拆分虚拟节点可保持图结构一致。
  \item 巷道横截面为宽 $4\,\mathrm{m}$、高 $3\,\mathrm{m}$ 的矩形，水流前锋对应水深 $0.1\,\mathrm{m}$。解释：在缺少糙率和局部阻力参数时，采用守恒的体积阈值近似，不求完整明渠水动力方程；矩形断面的流量--面积关系参考明渠水力学基本原理\cite{chow1959}。
  \item 高程不同的巷道只允许水从高端流向低端；水平巷道在每个水源传播层内按名义到达势定向。解释：这样既保留重力方向，又避免水平环路中水量无限循环。
  \item 分叉处对仍能容纳或传递水量的候选巷道平均分流；后到水源若存在不与既有来流对冲的出口，优先向这些出口传播。解释：附件未给出粗糙度、压差等加权依据，平均分流是参数最少且可复核的基准近似。
  \item 各突水点自题面给定时刻起以 $30\,\mathrm{m^3/min}$ 恒定供水；多水源共享巷道蓄水容量，巷道饱和后的来水记为边界积水。解释：该处理保证系统总水量守恒。
  \item 矿工接到通知后立即按规划路线连续行进，不在节点等待；水深达到 $0.3\,\mathrm{m}$ 后禁止进入或继续穿越。解释：该约束与应急逃生场景一致，也使边到达函数满足先出发不晚到达的标号设定条件。
\end{enumerate}
